% Part: incompleteness
% Chapter: arithmetization-syntax
% Section: coding-terms

\documentclass[../../../include/open-logic-section]{subfiles}

\begin{document}

\olfileid{inc}{art}{trm}
\olsection{ترميز الحدود}

\begin{explain}
الحد متتالية رموز ذات بناء مخصوص، تنشأ استقرائيًا
من الثوابت والمتغيرات بقواعد تكوين الحدود.
ولما أمكن ترميز الرموز، ثم متتاليات أعدادها،
لم يشق تعيين عدد غودل لكل حد. وإنما المطلوب
الأدق أن نثبت أن تمييز عدد غودل لحد حسن
التكوين قابل للحساب، بل عودي بدائي.
\end{explain}

!!^{variable}s وال!!{constant}s أول الحدود وأبسطها.
واختبار كون ${\OLId{latin-ordinary}{x}}$ عدد غودل لأحدها يسير:
فـ$\fn{Var}({\OLId{latin-ordinary}{x}})$ تصدق إذا كان ${\OLId{latin-ordinary}{x}}$ هو
$\Gn{\Obj v_{\OLId{latin-ordinary}{i}}}$ لبعض~${\OLId{latin-ordinary}{i}}$. أي إن ${\OLId{latin-ordinary}{x}}$
يرمز متتالية طولها~${\OLMathNumeral{1}}$، وعنصرها الوحيد
$({\OLId{latin-ordinary}{x}})_{\OLMathNumeral{0}}$ رمز !!{variable}~$\Obj v_{\OLId{latin-ordinary}{i}}$.
فتكون قيمة ${\OLId{latin-ordinary}{x}}$ هي $\tuple{\tuple{{\OLMathNumeral{1}}, {\OLId{latin-ordinary}{i}}}}$
لبعض~${\OLId{latin-ordinary}{i}}$. ونظير ذلك $\fn{Const}({\OLId{latin-ordinary}{x}})$:
تصدق إذا كان ${\OLId{latin-ordinary}{x}}$ هو $\Gn{\Obj c_{\OLId{latin-ordinary}{i}}}$
لبعض~${\OLId{latin-ordinary}{i}}$. والعلاقتان عوديتان بدائيتان،
إذ الشاهد ${\OLId{latin-ordinary}{i}}$ إن وجد كان $< {\OLId{latin-ordinary}{x}}$:
\begin{align*}
  \fn{Var}({\OLId{latin-ordinary}{x}}) & \defiff \bexists{{\OLId{latin-ordinary}{i}}<{\OLId{latin-ordinary}{x}}}{{\OLId{latin-ordinary}{x}} = \tuple{\tuple{{\OLMathNumeral{1}}, {\OLId{latin-ordinary}{i}}}}}\\
  \fn{Const}({\OLId{latin-ordinary}{x}}) & \defiff \bexists{{\OLId{latin-ordinary}{i}}<{\OLId{latin-ordinary}{x}}}{{\OLId{latin-ordinary}{x}} = \tuple{\tuple{{\OLMathNumeral{2}}, {\OLId{latin-ordinary}{i}}}}}
\end{align*}

\begin{prop}
\ollabel{prop:term-primrec}
كل من $\fn{Term}({\OLId{latin-ordinary}{x}})$ و$\fn{ClTerm}({\OLId{latin-ordinary}{x}})$ عودية بدائية،
ومعناهما أن ${\OLId{latin-ordinary}{x}}$ عدد غودل لحد، ولحد مغلق، على الترتيب.
\end{prop}

\begin{proof}
متتالية الرموز~${\OLId{latin-ordinary}{s}}$ حد إذا وفقط إذا وجدت
متتالية حدود ${\OLId{latin-ordinary}{s}}_{\OLMathNumeral{0}}$، \dots، ${\OLId{latin-ordinary}{s}}_{{\OLId{latin-ordinary}{k}}-{\OLMathNumeral{1}}}={\OLId{latin-ordinary}{s}}$
تسجل بناء~${\OLId{latin-ordinary}{s}}$ من ال!!{constant}s
وال!!{variable}s بقواعد التكوين.
ولصحة متتالية البناء لا بد، لكل ${\OLId{latin-ordinary}{i}}<{\OLId{latin-ordinary}{k}}$،
من أحد الأمور الآتية:
\begin{enumerate}
\item ${\OLId{latin-ordinary}{s}}_{\OLId{latin-ordinary}{i}}$ !!a{variable}~$\Obj v_{\OLId{latin-ordinary}{j}}$؛
\item أو ${\OLId{latin-ordinary}{s}}_{\OLId{latin-ordinary}{i}}$ !!a{constant}~$\Obj c_{\OLId{latin-ordinary}{j}}$؛
\item أو ${\OLId{latin-ordinary}{s}}_{\OLId{latin-ordinary}{i}}$ مؤلف من ${\OLId{latin-ordinary}{n}}$ حدود
  ${\OLId{latin-ordinary}{t}}_{\OLMathNumeral{1}}$، \dots، ${\OLId{latin-ordinary}{t}}_{\OLId{latin-ordinary}{n}}$ سبقت الموضع~${\OLId{latin-ordinary}{i}}$،
  باستعمال !!{function}~$\Obj f^{\OLId{latin-ordinary}{n}}_{\OLId{latin-ordinary}{j}}$ ذات ${\OLId{latin-ordinary}{n}}$ مواضع.
\end{enumerate}
فنحتاج إلى التعبير عن هذا الشرط على أعداد غودل
بدوال وعلاقات عودية بدائية وكمّات محدودة؛
ليثبت أن العلاقة المناظرة عودية بدائية.

ليكن ${\OLId{latin-ordinary}{y}}$ رمز المتتالية ${\OLId{latin-ordinary}{s}}_{\OLMathNumeral{0}}$، \dots، ${\OLId{latin-ordinary}{s}}_{{\OLId{latin-ordinary}{k}}-{\OLMathNumeral{1}}}$،
أي ${\OLId{latin-ordinary}{y}}=\tuple{\Gn{{\OLId{latin-ordinary}{s}}_{\OLMathNumeral{0}}}, \dots, \Gn{{\OLId{latin-ordinary}{s}}_{{\OLId{latin-ordinary}{k}}-{\OLMathNumeral{1}}}}}$.
وكونه رمز متتالية تكوين للحد ذي عدد غودل~${\OLId{latin-ordinary}{x}}$
يتوقف على تحقق أحد الشروط الآتية لكل ${\OLId{latin-ordinary}{i}}<{\OLId{latin-ordinary}{k}}$:
\begin{enumerate}
\item $\fn{Var}(({\OLId{latin-ordinary}{y}})_{\OLId{latin-ordinary}{i}})$؛ أو
\item $\fn{Const}(({\OLId{latin-ordinary}{y}})_{\OLId{latin-ordinary}{i}})$؛ أو
\item ثمة ${\OLId{latin-ordinary}{n}}$ وعدد~${\OLId{latin-ordinary}{z}}=\tuple{{\OLId{latin-ordinary}{z}}_{\OLMathNumeral{1}}, \dots, {\OLId{latin-ordinary}{z}}_{\OLId{latin-ordinary}{n}}}$ على أن يكون كل ${\OLId{latin-ordinary}{z}}_{\OLId{latin-ordinary}{l}}$
  مساويًا لقيمة $({\OLId{latin-ordinary}{y}})_{{\OLId{latin-ordinary}{i}}'}$ من أجل ${\OLId{latin-ordinary}{i}}'<{\OLId{latin-ordinary}{i}}$، و
\[
({\OLId{latin-ordinary}{y}})_{\OLId{latin-ordinary}{i}} = \Gn{\Obj f^{\OLId{latin-ordinary}{n}}_{\OLId{latin-ordinary}{j}}(} \concat \fn{flatten}({\OLId{latin-ordinary}{z}}) \concat \Gn{)},
\]
\end{enumerate}
مع اشتراط $({\OLId{latin-ordinary}{y}})_{{\OLId{latin-ordinary}{k}}-{\OLMathNumeral{1}}}={\OLId{latin-ordinary}{x}}$. (والدالة $\fn{flatten}({\OLId{latin-ordinary}{z}})$ تأخذ المتتالية
$\tuple{\Gn{{\OLId{latin-ordinary}{t}}_{\OLMathNumeral{1}}}, \dots, \Gn{{\OLId{latin-ordinary}{t}}_{\OLId{latin-ordinary}{n}}}}$ إلى $\Gn{{\OLId{latin-ordinary}{t}}_{\OLMathNumeral{1}}, \dots, {\OLId{latin-ordinary}{t}}_{\OLId{latin-ordinary}{n}}}$، وتكون
عودية بدائية.)

وفي الشرط الثالث، كل من ${\OLId{latin-ordinary}{j}}$ و${\OLId{latin-ordinary}{n}}$، وأعداد
غودل ${\OLId{latin-ordinary}{z}}_{\OLId{latin-ordinary}{l}}$ للحدود ${\OLId{latin-ordinary}{t}}_{\OLId{latin-ordinary}{l}}$، ورمز المتتالية
${\OLId{latin-ordinary}{z}}$، أي $\tuple{{\OLId{latin-ordinary}{z}}_{\OLMathNumeral{1}}, \dots, {\OLId{latin-ordinary}{z}}_{\OLId{latin-ordinary}{n}}}$،
دون~${\OLId{latin-ordinary}{y}}$. ونستبدل ${\OLId{latin-ordinary}{k}}$ بـ$\len{{\OLId{latin-ordinary}{y}}}$.
فبذلك تصاغ عبارة «${\OLId{latin-ordinary}{y}}$ رمز متتالية تكوين
لحد عدد غودل له~${\OLId{latin-ordinary}{x}}$» بدوال وعلاقات
عودية بدائية وكمّات محدودة.

ويبقى حصر ${\OLId{latin-ordinary}{y}}$ بحد عودي بدائي. فإذا كان
${\OLId{latin-ordinary}{x}}$ عدد غودل لحد، أمكن بناؤه بمتتالية
لا يزيد عدد حدودها على $\len{{\OLId{latin-ordinary}{x}}}$؛
إذ كل حد في متتالية تكوين~${\OLId{latin-ordinary}{s}}$ يبدأ
في موضع من~${\OLId{latin-ordinary}{s}}$، ولا يبدأ حدان جزئيان
في موضع واحد. وعدد غودل لكل حد جزئي
من~${\OLId{latin-ordinary}{s}}$ هو $ \le {\OLId{latin-ordinary}{x}}$. فتوجد متتالية
تكوين رمزها $\le {\OLId{latin-ordinary}{p}}_{{\OLId{latin-ordinary}{k}}-{\OLMathNumeral{1}}}^{{\OLId{latin-ordinary}{k}}({\OLId{latin-ordinary}{x}}+{\OLMathNumeral{1}})}$،
حيث ${\OLId{latin-ordinary}{k}}=\len{{\OLId{latin-ordinary}{x}}}$.

وأما $\fn{ClTerm}$، فنحصل على حجة عوديتها
البدائية بحذف حالة ال!!{variable}s من البناء.
\end{proof}

\begin{prob}
برهن عودية الدالة $\fn{flatten}({\OLId{latin-ordinary}{z}})$ البدائية، وهي التي ترسل المتتالية
$\tuple{\Gn{{\OLId{latin-ordinary}{t}}_{\OLMathNumeral{1}}}, \dots, \Gn{{\OLId{latin-ordinary}{t}}_{\OLId{latin-ordinary}{n}}}}$ إلى $\Gn{{\OLId{latin-ordinary}{t}}_{\OLMathNumeral{1}}, \dots, {\OLId{latin-ordinary}{t}}_{\OLId{latin-ordinary}{n}}}$.
\end{prob}

\begin{prop}\ollabel{prop:num-primrec}
  تحسب الدالة $\fn{num}({\OLId{latin-ordinary}{n}})=\Gn{\num{{\OLId{latin-ordinary}{n}}}}$ بالعودية البدائية.
\end{prop}

\begin{proof}
  يكفي التعريف الآتي لـ$\fn{num}({\OLId{latin-ordinary}{n}})$ بالعودية البدائية:
  \begin{align*}
    \fn{num}({\OLMathNumeral{0}}) & = \Gn{\Obj 0}\\
    \fn{num}({\OLId{latin-ordinary}{n}}+{\OLMathNumeral{1}}) & = \Gn{\prime(} \concat \fn{num}({\OLId{latin-ordinary}{n}}) \concat \Gn{)}.
  \end{align*}
\end{proof}

\end{document}
